Let $E$ be the selected element width in bytes, $N=VLEN/E$, and $G$ be
the typed governing-predicate image whose significant bits are
\texttt{Pg[iE]} for $0\leq i<N$.  \texttt{Pg} and \texttt{Ps} shall name
different (:term:VECTOR.terms.predicate_register|plural:).  Before attempting a lane, VSCATTER replaces
\texttt{Ps} by its typed image \texttt{Ps AND G}; all nonsignificant bits are
cleared.

If \texttt{G AND NOT Ps} is empty, the instruction completes normally with
\texttt{Ps == G}.  Otherwise it selects one pending lane $i$, calculates that
lane's address as specified by the resumable gather/scatter rules in
Section~\ref{section:vector-register-model}, and issues exactly one $E$-byte
normal-memory store of source lane $i$ through the default data segment.  On
success it sets \texttt{Ps[iE]} and commits the store and completion bit at a
lane-completion boundary.  If more lanes remain pending, the PC continues to
identify this VSCATTER and execution may admit an asynchronous event before
another lane is attempted.  A lane-completion boundary is not an instruction
retirement or a debug-trace boundary.

Lane selection order is not architecturally specified.  When stores from two
lanes overlap in physical memory, the final value of each overlapping byte is
therefore unspecified; every successful lane nevertheless performs exactly
one $E$-byte store and sets its completion bit.  A failure reported before the
store's point of no return leaves that bit clear and is retry-safe.  A failure
reported after the store became irrevocable sets the bit and is reported later
as an imprecise, non-retry-safe machine check rather than as a restartable
synchronous fault.  Software emulation shall perform or explicitly waive the
store before setting the completion bit.  The address and source vector may
name the same (:term:VECTOR.terms.vector_register:).  VSCATTER preserves FLAGS and is not permitted for
MMIO.
